%!TEX TS-program = pdflatex
\newcommand{\printpraesenzlsg}{false}
\newcommand{\printloesungen}{true}
\newcommand{\printbewertungen}{false}
\newcommand{\blattnummer}{4}
\documentclass[%
  paper=a4,
  fontsize=10pt,
  ngerman
  ]{scrartcl}

% Basics für Codierung und Sprache
% ===========================================================
  \usepackage{scrtime}
  \usepackage{etex}
  \usepackage{shellesc}
  \usepackage[final]{graphicx}
  \usepackage[utf8]{inputenc}
  \usepackage{babel}
  \usepackage[german=quotes]{csquotes}
  \usepackage{datetime2}
  \usepackage{ifthen}
  \usepackage{dingbat}
% ===========================================================

%% Fonts und Typographie
%% ===========================================================
  \usepackage{charter}
  \usepackage{nimbusmononarrow}
  \usepackage[babel=true,final,tracking=smallcaps]{microtype}
  \DisableLigatures{encoding = T1, family = tt* }
  \usepackage{ellipsis}
  \usepackage{listings}
  
% ===========================================================

% Farben
% ===========================================================
  \usepackage[usenames,x11names,table]{xcolor}
  \definecolor{fbblau}{HTML}{3078AB}
  \definecolor{blackberry}{rgb}{0.53, 0.0, 0.25}
% ===========================================================

% Mathe-Pakete und -Einstellungen
% ===========================================================
  \usepackage{mathtools}
  \usepackage{amssymb}
  \usepackage[bigdelims]{newtxmath}
  \allowdisplaybreaks
  \usepackage{bm}
  \usepackage{wasysym}
% ===========================================================

% TikZ
% ===========================================================
  \usepackage{tikz}
  \usepackage{tikz-cd}
  \usepackage{pgf}
  \usetikzlibrary{matrix}
  \usetikzlibrary{positioning}
  \usetikzlibrary{arrows}
  \usetikzlibrary{shapes}
  \tikzset{>=Latex}
  \usepackage[mode=image]{standalone}
% ===========================================================

% Seitenlayout, Kopf-/Fußzeile
% ===========================================================
  \usepackage{scrlayer-scrpage}
  \usepackage[top=3cm, bottom=2.5cm, left=2.5cm, right=2cm]{geometry}  
  \pagestyle{scrheadings}
  \clearscrheadfoot 
  \setheadsepline{0.4pt}
  \setfootsepline{0.4pt}
  \setkomafont{pagehead}{\normalfont\footnotesize}
  \setkomafont{pagefoot}{\normalfont\footnotesize}
  \lohead[]{Übung zur Vorlesung \textit{Embedded Systems Engineering} -- Wintersemester 2026}
  \rohead[]{%
    \ifthenelse{\boolean{loesungen}}{\textbf{Lösungen zu Blatt \blattnummer}}{Blatt \blattnummer}}
  %\rofoot[Version: \DTMnow]{%
  %  \ifthenelse{\boolean{loesungen}}{Version: \DTMnow}{}}
  %\lofoot[\jobname.tex]{%
  %  \ifthenelse{\boolean{loesungen}}{\jobname.tex}{}}
  \cfoot[\thepage]{\thepage}
  \raggedbottom
  \usepackage{setspace}          
  \linespread{1.2}
  \setlength{\parindent}{0pt}
  \setlength{\parskip}{0.5\baselineskip}
  \usepackage[all]{nowidow}
  \usepackage{lscape}
% ===========================================================

% Hyperref
% ===========================================================
  \usepackage[%
    hidelinks,
    pdfpagelabels,
    bookmarksopen=true,
    bookmarksnumbered=true,
    linkcolor=black,
    urlcolor=fbblau,
    plainpages=false,
    pagebackref,
    citecolor=black,
    hypertexnames=true,
    pdfborderstyle={/S/U},
    linkbordercolor=fbblau,
    colorlinks=false,
    backref=false]{hyperref}
  \hypersetup{final}
  \makeatletter
  \g@addto@macro{\UrlBreaks}{\UrlOrds}
  \makeatother
% ===========================================================

% Listen und Tabellen
% ===========================================================
  \usepackage{multicol}
  \usepackage{multirow}
  \usepackage[shortlabels]{enumitem}
  \setlist[enumerate]{font=\bfseries,itemsep=0.25\baselineskip}
  \setlist[itemize]{label=$\triangleright$,itemsep=-0.25\baselineskip}
  \usepackage{tabularx}
  \usepackage{array}
  \usepackage{listings}
  \lstset{breaklines=true,
    basicstyle=\ttfamily,
    frame=tblr,
    language=c,
    numbers=left}
  \usepackage{include/mips}

\definecolor{lightgreen}{rgb}{0,0.7,0}
\definecolor{gray}{rgb}{0.5,0.5,0.5}
\definecolor{mauve}{rgb}{0.58,0,0.82}

\lstdefinestyle{mips}{ %
  language=[mips]Assembler,       % the language of the code
  keywordstyle=\color{blue},          % keyword style
  keywordstyle=[2]\color{magenta},     % keyword style
  keywordstyle=[3]\color{red},          % keyword style
  commentstyle=\color{lightgreen},       % comment style
  stringstyle=\color{mauve},         % string literal style
  escapeinside={\%*}{*)},            % if you want to add a comment within your code
  morekeywords={*,...}               % if you want to add more keywords to the set
}

\newcommand{\mipsinline}[1]{\lstinline[style=mips,columns=fixed]{#1}}
\newcommand{\mipsregister}[1]{\mipsinline{\$#1}}
  \newcommand{\qed}{\hfill \ensuremath{\Box}}
  \newcommand{\cmd}[1]{\texttt{#1}}
% ===========================================================

% Sonstiges
% ===========================================================
  \usepackage[textsize=small,color=Red1!80!OrangeRed1!80]{todonotes}
  \usepackage{lipsum}
% ===========================================================

% Shortcuts and Math Commands
% ===========================================================
  \newcommand{\BB}{\mathbb{B}}      % Bool
  \newcommand{\CC}{\mathbb{C}}      % komplexe Zahlen
  \newcommand{\NN}{\mathbb{N}}      % nat. Zahlen
  \newcommand{\QQ}{\mathbb{Q}}      % rat. Zahlen
  \newcommand{\RR}{\mathbb{R}}      % reelle Zahlen
  \newcommand{\ZZ}{\mathbb{Z}}      % ganze Zahlen
  \newcommand{\oh}{\mathcal{O}}      % Landau-O
  \newcommand{\ol}[1]{\overline{#1}}    % überstreichen
  \newcommand{\wt}[1]{\widetilde{#1}}    % überschlängeln
  \newcommand{\wh}[1]{\widehat{#1}}    % überdachen
  \newcommand{\entspricht}{\mathrel{\widehat{=}}}
  
  \DeclarePairedDelimiter{\absolut}{\lvert}{\rvert}    % Betrag
  \DeclarePairedDelimiter{\ceiling}{\lceil}{\rceil}    % aufrunden
  \DeclarePairedDelimiter{\Floor}{\lfloor}{\rfloor}    % aufrunden
  \DeclarePairedDelimiter{\Norm}{\lVert}{\rVert}      % Norm
  \DeclarePairedDelimiter{\sprod}{\langle}{\rangle}    % spitze Klammern
  \DeclarePairedDelimiter{\enbrace}{(}{)}          % runde Klammern
  \DeclarePairedDelimiter{\benbrace}{\lbrack}{\rbrack}  % eckige Klammern
  \DeclarePairedDelimiter{\penbrace}{\{}{\}}        % geschweifte Klammern
  \newcommand{\Underbrace}[2]{{\underbrace{#1}_{#2}}}   % bessere Unterklammerungen
  
  % Kurzschreibweisen für Faule und Code-Vervollständigung
  \newcommand{\abs}[1]{\absolut*{#1}}
  \newcommand{\ceil}[1]{\ceiling*{#1}}
  \newcommand{\flo}[1]{\Floor*{#1}}
  \newcommand{\no}[1]{\Norm*{#1}}
  \newcommand{\sk}[1]{\sprod*{#1}}
  \newcommand{\enb}[1]{\enbrace*{#1}}
  \newcommand{\penb}[1]{\penbrace*{#1}}
  \newcommand{\benb}[1]{\benbrace*{#1}}
  \newcommand{\stack}[2]{\makebox[1cm][c]{$\stackrel{#1}{#2}$}}

% Übungszettel-Krams
% ===========================================================
  \usepackage[tikz]{mdframed}
  \newmdenv[%
    backgroundcolor = gray!20,
    linewidth=0pt,
    skipabove = 0.5cm,
    skipbelow = 0cm
  ]{notes}
  
  % Lösungsausgabe
  \newboolean{praesenzlsgn}
  \setboolean{praesenzlsgn}{\printpraesenzlsg}
  \newboolean{loesungen}
  \setboolean{loesungen}{\printloesungen}
  % Bewertungsschema-Ausgabe
  \newboolean{bewertungen}
  \setboolean{bewertungen}{\printbewertungen}

  \newcommand{\iforiginal}[1]{\ifthenelse{\boolean{praesenzlsgn}}{}{#1}}
  \newcommand{\ifpraesenzlsg}[1]{\ifthenelse{\boolean{praesenzlsgn}}{#1}{}}
  \newcommand{\ifloesung}[1]{\ifthenelse{\boolean{loesungen}}{#1}{}}
  \newcommand{\ifbewertung}[1]{\ifthenelse{\boolean{bewertungen}}{#1}{}}

  \newcommand{\umbruchoriginal}{\iforiginal{\newpage}}        % Seitenumbruch nur im Original
  \newcommand{\umbruchpraesenzlsg}{\ifpraesenzlsg{\newpage}}  % Seitenumbruch nur in Lösung d. Präsenzaufgabe
  \newcommand{\umbruchlsg}{\ifloesung{\newpage}}              % Seitenumbruch nur in Lösung
  \newcommand{\umbruchbewertung}{\ifbewertung{\newpage}}      % Seitenumbruch nur in Lösung mit Bewertungsschema
  \newcommand{\umbruchnurpraesenz}{\ifthenelse{\boolean{loesungen}}{}{\umbruchpraesenzlsg}}
  \newcommand{\umbruchnurlsg}{\ifthenelse{\boolean{bewertungen}}{}{\umbruchlsg}}

  \usepackage{comment}
  \newenvironment{praesenzlsg}%
    {\ifthenelse{\boolean{praesenzlsgn}}{\textbf{--- Präsenzaufgabe Anfang ---} \mbox{} \newline}{}}%
	{\ifthenelse{\boolean{praesenzlsgn}}{\mbox{} \newline \textbf{--- Präsenzaufgabe Ende ---}}{}}
  \newenvironment{loesung}%
    {\ifthenelse{\boolean{loesungen}}{\textbf{--- Lösung Anfang ---} \mbox{} \newline}{}}%
    {\ifthenelse{\boolean{loesungen}}{\mbox{} \newline \textbf{--- Lösung Ende ---} \ifthenelse{\boolean{bewertungen}}{}{\umbruchlsg}}{}}
  \newenvironment{bewertung}%
    {\ifthenelse{\boolean{bewertungen}}{\textbf{--- Bewertungsschema Anfang ---} \mbox{} \newline}{}}%
    {\ifthenelse{\boolean{bewertungen}}{\mbox{} \newline \textbf{--- Bewertungsschema Ende ---} \umbruchlsg}{}}

  \ifthenelse{\boolean{praesenzlsgn}}{}{\excludecomment{praesenzlsg}}
  \ifthenelse{\boolean{loesungen}}{}{\excludecomment{loesung}}
  \ifthenelse{\boolean{bewertungen}}{}{\excludecomment{bewertung}}

  % Aufgabennummerierung
  \newcounter{aufgabennummer} 
  \setcounter{aufgabennummer}{1}
  \newcommand{\aufgabe}[1]{\textbf{Aufgabe \blattnummer.\theaufgabennummer} \stepcounter{aufgabennummer} \hfill (#1 Punkte)}
  %\newcommand{\aufgabetitel}[2]{\iforiginal{\vspace{0.5cm}} \textbf{Aufgabe \blattnummer.\theaufgabennummer \ (#2)} \stepcounter{aufgabennummer} \hfill (#1 Punkte)}
  \newcommand{\aufgabetitel}[2]{\iforiginal{\vspace{0.5cm}} \textbf{Aufgabe \theaufgabennummer \ (#2)} \stepcounter{aufgabennummer} \hfill (#1 Punkte)}
  \newcommand{\praesenztitel}[1]{\iforiginal{\vspace{0.5cm}} \textbf{Aufgabe \blattnummer.\theaufgabennummer \ (#1)} \stepcounter{aufgabennummer} \hfill \smallpencil}

\newcommand{\danger}{\raisebox{0.5mm}{\fontencoding{U}\fontfamily{futs}\selectfont\char 66\relax}}
\newcommand{\ddanger}[1]{{\danger} \textbf{#1} {\danger}}

% Blatt-Thema (optional)
\newcommand{\blatthema}{} % default: leer
\newcommand{\setblatthema}[1]{\renewcommand{\blatthema}{#1}}

% Rechnerstrukturen-Krams
% ===========================================================
  \newcommand{\Band}{\ensuremath{\mathbin{\texttt{AND}}}}
  \newcommand{\Bor}{\ensuremath{\mathbin{\texttt{OR}}}}
  \newcommand{\Bnot}{\ensuremath{\mathbin{\texttt{NOT}}}}
  \newcommand{\Bnand}{\ensuremath{\mathbin{\texttt{NAND}}}}
  \newcommand{\Bnor}{\ensuremath{\mathbin{\texttt{NOR}}}}
  \newcommand{\Bxor}{\ensuremath{\mathbin{\texttt{XOR}}}}
  \newcommand{\Bxnor}{\ensuremath{\mathbin{\texttt{XNOR}}}}
  \newcommand{\Bone}{\ensuremath{\mathbin{\texttt{ONE}}}}
  \newcommand{\Bzero}{\ensuremath{\mathbin{\texttt{ZERO}}}}
  \newcommand{\Binary}[1]{\ensuremath{\texttt{#1}}\xspace}
  \newcommand{\ieee}[3]{\fbox{\ensuremath{\texttt{#1}}}\fbox{\ensuremath{\texttt{#2}}}\fbox{\ensuremath{\texttt{#3}}}}
  \DeclareMathOperator{\DNF}{DNF}
  \DeclareMathOperator{\KNF}{KNF}
  \DeclareMathOperator{\xor}{\oplus}
  \newcommand{\dc}{\star} %don't-care-Wert


% Create Code Listings
% ===========================================================
  \lstset{
  basicstyle=\ttfamily,
  numbers=left,
  numberstyle=\small,
  stepnumber=1,
  numbersep=10pt,
  frame=single,
  columns=fullflexible,
  keepspaces=true,
  showstringspaces=false,
  xleftmargin=1.5em
}
 
\usetikzlibrary{automata,arrows.meta,positioning,fit,backgrounds}
\tikzset{zustand/.style={draw,rounded corners,align=center,minimum height=9mm,minimum width=20mm},every edge/.style={draw,->},>=Stealth}
\lstset{basicstyle=\ttfamily\small,numbers=none,language={},xleftmargin=0pt}
\begin{document}
\section*{Lösungen zu Blatt 4: State Charts}
Die Definitionen aus dem Aufgabenblatt werden unverändert übernommen. Alle Modelle werden mit \texttt{@EventDriven} ausgeführt. Kantenbeschriftungen verwenden \texttt{Trigger [Guard] / Effekt}; Initial- und History-Entries besitzen keinen Trigger. Nicht behandelte Eingabeereignisse werden ignoriert. Die Diagramme und Tabellen beschreiben die im CREATE-Modellbereich anzulegenden Zustände und Transitionen.

\subsection*{Aufgabe 1: Fahrstuhl}
Vier Zustände unterscheiden Stillstand und Fahrt: \textit{EG}, \textit{Aufwärts}, \textit{OG}, \textit{Abwärts}. Die Fahrzeit wird durch Zeitevents modelliert, ohne zusätzliche Events oder Variablen.
\begin{center}
\begin{tikzpicture}[auto,font=\small]
\node[zustand] (eg) at (0,0) {EG};
\node[circle,fill,inner sep=2pt] (init) at (-3,0) {};
\node[zustand] (up) at (7,0) {Aufwärts};
\node[zustand] (og) at (7,-3) {OG};
\node[zustand] (down) at (0,-3) {Abwärts};
\draw[->] (init) -- node[above,align=center]{/ location = 2} (eg);
\path (eg) edge node{Button.up} (up)
(up) edge node[right,align=left]{after 3s /\\location = 1;\\raise Status.arrived} (og)
(og) edge node[below]{Button.down} (down)
(down) edge node[left,align=right]{after 3s /\\location = 2;\\raise Status.arrived} (eg);
\end{tikzpicture}
\end{center}
\textbf{Exakte Transitionen:}
\begin{lstlisting}
Initial -> EG: / location = 2
EG -> Aufwaerts: Button.up
Aufwaerts -> OG:
    after 3s / location = 1; raise Status.arrived
OG -> Abwaerts: Button.down
Abwaerts -> EG:
    after 3s / location = 2; raise Status.arrived
\end{lstlisting}
Die Schreibweise \texttt{Quelle -> Ziel:} dient hier und im Folgenden nur der Dokumentation; in CREATE wird lediglich der Text hinter dem Doppelpunkt an die jeweilige Kante geschrieben.

Während der Fahrt werden Tastenereignisse ignoriert. Insbesondere darf kein Selbstübergang in einem Fahrzustand den 3-Sekunden-Timer neu starten. \texttt{location} bleibt bis zur Ankunft das zuletzt erreichte Stockwerk; ein zusätzlicher Wert für „unterwegs“ ist nicht nötig. Beim Initialisieren wird kein Ankunftsevent ausgegeben.

\textbf{Prüffolge:} Start: EG, \texttt{location = 2}; \texttt{Button.up}: Aufwärts; nach 3 Sekunden: OG, \texttt{location = 1} und einmal \texttt{Status.arrived}. Die Rückfahrt verläuft symmetrisch.
\newpage
\subsection*{Aufgabe 2a: Reinigungsroboter}
Oberste Ebene: \textit{Aus} (initial) und der zusammengesetzte Zustand \textit{Ein}. \textit{Ein} besitzt zwei orthogonale Regionen: \textit{Betriebsmodus} und \textit{Hindernisüberwachung}. Beim Einschalten wird über einen \textbf{Shallow-History-Entry} in der Modusregion der letzte Modus reaktiviert; ohne gespeicherte History führt dessen Default-Kante nach \textit{Warten}.
\begin{center}
\begin{tikzpicture}[font=\small,auto]
\draw[rounded corners] (-1,-1) rectangle (11,-6.8);
\node[anchor=north west,font=\bfseries] at (-.9,-1.05) {Ein};
\draw[dashed] (-1,-4.8) -- (11,-4.8);
\node[anchor=north west] at (-.9,-1.55) {Region Betriebsmodus};
\node[anchor=north west] at (-.9,-4.85) {Region Hindernisüberwachung};
\node[zustand,initial,initial text={}] (aus) at (0,1) {Aus};
\node[circle,draw] (h) at (0,-2.8) {H};
\node[zustand] (w) at (2.4,-2.8) {Warten};
\node[zustand] (s) at (5.8,-2.8) {Saugen};
\node[zustand] (m) at (9,-2.8) {Wischen};
\draw[->] (aus) -- node[right,align=left,pos=.28]{Remote.power /\\raise Status.ein} (h);
\draw[->] (10,-1) -- (10,1) -- node[above]{Remote.power / raise Status.aus} (aus);
\draw[->] (h) -- (w);
\path (w) edge node[above]{T1} (s)
(s) edge node[above]{T2} (m)
(m) edge[bend left=30] node[below]{T3} (w)
(s) edge[bend left=45] node[below]{T4} (w);
\node[circle,fill,inner sep=2pt] (i) at (0,-6) {};
\node[zustand] (sensor) at (4,-6) {Überwachen};
\draw[->] (i) -- (sensor);
\path (sensor) edge[loop right] node[right,align=left]{Hindernis.detect /\\raise Status.hindernis} (sensor);
\end{tikzpicture}
\end{center}
Die Moduskanten sind aus Platzgründen abgekürzt; verbindlich sind folgende vollständigen Beschriftungen:
\begin{lstlisting}
T1: Warten -> Saugen: Remote.mode / raise Status.saugen
T2: Saugen -> Wischen: Remote.mode / raise Status.wischen
T3: Wischen -> Warten:
    Remote.mode, Remote.stop / raise Status.warten
T4: Saugen -> Warten: Remote.stop / raise Status.warten
\end{lstlisting}
\textbf{History und Hierarchie.} Die Einschaltkante zielt auf den History-Entry der Modusregion in \textit{Ein}; die zweite Region wird dabei über ihren normalen Initial-Entry betreten. Die Ausschaltkante beginnt am Rand des gesamten Zustands \textit{Ein}. Sie beendet beide Regionen und bewahrt den letzten Modus. Shallow History reicht, da die drei Modi keine Unterzustände haben.

\textbf{Ausgaben.} Statusmeldungen für Moduswechsel stehen auf den Kanten, nicht als Entry-Aktionen. Dadurch wird beim Wiederherstellen nur \texttt{Status.ein} ausgegeben. Optional kann die Default-Kante $H\to Warten$ den Effekt \texttt{/ raise Status.warten} tragen, um den Modus beim allerersten Start anzuzeigen. Der Hindernis-Selbstübergang betrifft nur \textit{Überwachen}, setzt also den Betriebsmodus nicht zurück. Im Zustand \textit{Aus} werden Sensor-, Modus- und Stoppereignisse ignoriert.
\newpage
\subsection*{Aufgabe 2b: Überwachung des Beutelfüllstands}
Die Zustände \textit{Messung}, \textit{Pause5}, \textit{Pause1} und \textit{Voll} genügen auf oberster Ebene. \textit{Voll} enthält die Unterzustände \textit{Frist} und \textit{Erinnert}. Jede Messanforderung führt nach \textit{Messung}; dessen Entry-Aktion lautet \texttt{entry / raise Sensor.check\_dust}.
\begin{center}
\begin{tikzpicture}[font=\small,auto,scale=.88,transform shape]
\node[zustand,initial,initial text={}] (mess) at (0,0) {Messung};
\node[zustand] (p5) at (8,1) {Pause5};
\node[zustand] (p1) at (8,-1.4) {Pause1};
\draw[rounded corners] (-1.5,-3.3) rectangle (10,-6.2);
\node[anchor=north west,font=\bfseries] at (-1.4,-3.35) {Voll};
\node[circle,fill,inner sep=2pt] (init) at (-.8,-5) {};
\node[zustand] (f) at (1,-5) {Frist};
\node[zustand] (e) at (7,-5) {Erinnert};
\draw[->] (init) -- (f);
\path (mess) edge[bend left=15] node[above]{Ergebnis $<90$} (p5)
(p5) edge[bend left=15] node[below]{after 5s} (mess)
(mess) edge[bend left=15] node[above]{Ergebnis $90\ldots98$ / Ton} (p1)
(p1) edge[bend left=15] node[below]{after 1s} (mess)
(f) edge node[above]{after 5s / raise Beutel.sound} (e);
\draw[->] (mess) -- node[right,align=left]{Ergebnis $\geq99$ / Ton} (0,-3.3);
\draw[->] (10,-4) -- (11,-4) -- (11,2.2) -- node[above]{Beutel.replaced} (0,2.2) -- (mess);
\end{tikzpicture}
\end{center}
\textbf{Exakte Modellbeschriftungen} (die Ergebniskanten im Diagramm sind abgekürzt):
\begin{lstlisting}
Initial -> Messung:                    // ohne Effekt
Messung: entry / raise Sensor.check_dust
Messung -> Pause5:
    Sensor.dust_checked [Sensor.dust_level < 90]
Messung -> Pause1:
    Sensor.dust_checked
    [Sensor.dust_level >= 90 && Sensor.dust_level < 99]
    / raise Beutel.sound
Messung -> Voll:
    Sensor.dust_checked [Sensor.dust_level >= 99]
    / raise Beutel.sound
Pause5 -> Messung: after 5s
Pause1 -> Messung: after 1s
Voll.Initial -> Frist:                 // ohne Effekt
Frist -> Erinnert: after 5s / raise Beutel.sound
Voll -> Messung: Beutel.replaced
\end{lstlisting}
\textbf{Warum zwei Unterzustände?} Bei Eintritt in \textit{Voll} ertönt der erste Ton. Nach 5 Sekunden folgt genau ein zweiter Ton; \textit{Erinnert} hat keinen Timer. \texttt{Beutel.replaced} verlässt aus beiden Unterzuständen den Verbundzustand und startet sofort eine neue Messung. Beim vorzeitigen Verlassen von \textit{Frist} wird dessen Timer beendet.

Die Guards sind disjunkt und für den zugesicherten Bereich 0--100 vollständig. Ausgewertet wird erst bei \texttt{Sensor.dust\_checked}, nie direkt nach der Messanforderung. Wartezeiten beginnen mit dem verarbeiteten Ergebnis. Nebenläufigkeit ist hier nicht nötig: Messung, Pause und Warten auf Austausch erfolgen nacheinander. Bei zeitgleichem Austausch und Timerablauf ist die Reihenfolge der Ereignisverarbeitung maßgeblich; die Spezifikation legt dafür keine zusätzliche Priorität fest.
\end{document}
